\section{Experimental Results}\label{sec:results}
\subsection{Published Experimental Baseline}\label{sec:baseline}
The conference experiments provide historical context for the
extensions~\cite{zone_secrypt_2026}. Their corpus contained 102 Windows
payloads: 38 x64 and 64 x86. Fifteen single and combined transformation
configurations were considered, with ten variants planned for each
input/configuration pair. Planned multiplicity is distinguished from
the number of outputs actually available.

Table~\ref{tab:historical} summarizes selected published aggregate
observations. The original-file malicious-verdict means were 21.67\%
for x64 and 22.63\% for x86. The published transformed-file averages were
substantially smaller, but functional outcomes varied across
configurations. These measurements belong to the original engine,
corpus, harness, and scan periods. They are not new experiments on the
extensions described above.

\begin{table}[t]
\caption{Selected published aggregate static-scan observations,
summarized from~\cite{zone_secrypt_2026}. Percentages describe the
published malicious-verdict fractions, not deployed endpoint detection.}
\label{tab:historical}
\centering\small
\begin{tabular}{@{}lrr@{}}
\toprule
Published condition & x64 & x86 \\
\midrule
Original payloads & 21.67\% & 22.63\% \\
Zone configurations, reported average & 0.57\% & 0.25\% \\
SGN comparison & 0.00\% & 10.22\% \\
\bottomrule
\end{tabular}
\end{table}

The conference comparison with SGN also considered memory dumps after
decoding~\cite{zone_secrypt_2026}. Such a submission tests recognition
of recovered bytes, whereas Zone's intended model retains a rewritten
representation. Neither observation directly measures the runtime
response of an installed endpoint product. The comparison is retained
as historical context rather than evidence about the new LLM or
wrapping paths.

The historical figures for size, entropy, and similarity characterized
representation changes. Their metrics remain distinct from functional
checks. The extended assessment uses the explicit definitions in
Section~\ref{sec:methodology}; any reanalysis of historical artifacts
must state its own grouping, deduplication, and representation choices.
\draftnote{The previous draft's detailed historical plots are omitted
pending reconciliation of category/value mappings and metric provenance.
The original sources are preserved in the backup. Finalize actual
historical output counts and identify any corrected reanalysis explicitly.}

\subsection{Code--Data Identification}\label{sec:data_results}
\draftnote{Insert the frozen identification campaign here. If no
independent reference annotations are available, present agreement and
case studies and retain that scope in the abstract and conclusions.}
\begin{table}[t]
\caption{Identification assessment. A dash denotes an unavailable
reference-based measure; pending denotes an experiment not yet finalized.}
\label{tab:data_results}
\centering\small
\begin{tabularx}{\linewidth}{@{}Yllll@{}}
\toprule
Method and output stage & Inputs & Precision$_D$ & Recall$_D$ & $F_{1,D}$ \\
\midrule
Symbolic/heuristic candidates & \pending & \pending & \pending & \pending \\
Symbolic/heuristic accepted & \pending & \pending & \pending & \pending \\
LLM candidates & \pending & \pending & \pending & \pending \\
LLM accepted & \pending & \pending & \pending & \pending \\
\bottomrule
\end{tabularx}
\end{table}
The completed analysis will distinguish interval agreement, uncertainty,
filtering, and downstream behavior. Examples will identify whether a
disagreement originates in the proposed data ranges, reference support,
or subsequent transformation. Identical final binaries will be reported
as pipeline agreement rather than used as a substitute for reference
annotations.

\subsection{Functional Assessment of the Extended Engine}
\label{sec:functional_results}
\draftnote{Insert preparation-only and native-transformation outcomes,
with the input-selection rule, architecture, snapshot, and harness.
Keep historical and current MessageBox oracles distinct.}
\begin{table}[t]
\caption{Functional assessment structure. Report counts before
percentages; $N$, $B$, and $W$ are defined in the methodology.}
\label{tab:functional_results}
\centering\small
\begin{tabularx}{\linewidth}{@{}Ylllll@{}}
\toprule
Condition & Arch. & $N$ & $B$ & $W/N$ & $W/B$ \\
\midrule
Data handling only & x64 & \pending & \pending & \pending & \pending \\
Data handling only & x86 & \pending & \pending & \pending & \pending \\
Native transformations & x64 & \pending & \pending & \pending & \pending \\
Native transformations & x86 & \pending & \pending & \pending & \pending \\
\bottomrule
\end{tabularx}
\end{table}
The failure analysis will distinguish construction refusal, incomplete
decoding, unsupported references, failed execution checks, and
uncompleted tests. Aggregate improvement will be interpreted alongside
changes in dependencies, implementation, and the observation protocol.

\subsection{Assembly Wrapping}\label{sec:wrapping_results}
\draftnote{Complete x64 outcomes from the frozen wrapping campaign.
Populate x86 rows only after its development version and evaluation are
fixed. Retain all planned cases, including unsuccessful builds.}
\begin{table}[t]
\caption{Wrapping assessment by architecture and configuration.
The x86 port remains development work in this draft.}
\label{tab:wrapping_results}
\centering\small
\begin{tabularx}{\linewidth}{@{}Yllll@{}}
\toprule
Condition & $N$ & $B/N$ & $W/N$ & $W/B$ \\
\midrule
x64 wrapping & \pending & \pending & \pending & \pending \\
x64 wrapping, combined & \pending & \pending & \pending & \pending \\
x86 wrapping & \multicolumn{4}{c}{\dev} \\
x86 wrapping, combined & \multicolumn{4}{c}{\dev} \\
\bottomrule
\end{tabularx}
\end{table}
The representation assessment will additionally report size, retained
regions, and unique output hashes. Any similarity score will identify
the decoding view. A result restricted to constructed outputs will be
accompanied by construction availability over all planned cases.

\subsection{LLM Transformation: Zero-Shot versus Two-Shot}
\label{sec:llm_results}
\draftnote{Insert the completed 0/2-shot campaign, preserving failed
and repeated runs. Report the number of distinct input payloads and
state whether functional verification participated in selection.}
\begin{table}[t]
\caption{LLM transformation outcomes. Initial and final acceptance
refer respectively to the first attempt and the full retry budget.}
\label{tab:llm_results}
\centering\small
\begin{tabularx}{\linewidth}{@{}Yllll@{}}
\toprule
Family / shots & Runs & Initial acc. & Final acc. & Functional \\
\midrule
EQ / 0 & \pending & \pending & \pending & \pending \\
EQ / 2 & \pending & \pending & \pending & \pending \\
NOP / 0 & \pending & \pending & \pending & \pending \\
NOP / 2 & \pending & \pending & \pending & \pending \\
BCF / 0 & \pending & \pending & \pending & \pending \\
BCF / 2 & \pending & \pending & \pending & \pending \\
IBP / 0 & \pending & \pending & \pending & \pending \\
IBP / 2 & \pending & \pending & \pending & \pending \\
\bottomrule
\end{tabularx}
\end{table}
The comparison will pair those outcomes with attempt counts,
sequence-based similarity, output uniqueness, elapsed time, and cost.
A smaller similarity value will not be interpreted as better functional
preservation, and acceptance after feedback will not be described as
success of the first zero-shot or two-shot answer.

\subsection{Static-Scan and Representation Measurements}
\draftnote{Insert completed scan records and matched representation
measurements. Pending scans are missing observations, never zero
detections. Keep the two verdict measures in distinct columns.}
\begin{table}[t]
\caption{Static-scan reporting structure. Completed hashes and group
membership are distinguished; $\overline d$ and $q$ have different definitions.}
\label{tab:scan_results}
\centering\small
\begin{tabularx}{\linewidth}{@{}Yllll@{}}
\toprule
Group & Unique hashes & Completed & $\overline d$ & $q$ \\
\midrule
Current original controls & \pending & \pending & \pending & \pending \\
Wrapping, by architecture & \pending & \pending & \pending & \pending \\
LLM transformation, by condition & \pending & \pending & \pending & \pending \\
LLM data handling, by condition & \pending & \pending & \pending & \pending \\
\bottomrule
\end{tabularx}
\end{table}
The final discussion will relate these observations to build and
functional outcomes on the same identified artifacts. Static verdicts
for a selected set of constructed binaries will not be generalized
to inputs for which construction or execution failed.
